\documentclass[11pt]{article}
\usepackage[margin=1in]{geometry}
\usepackage{amsmath,amssymb,amsthm}
\usepackage{mathtools}
\usepackage{hyperref}

\newtheorem{lemma}{Lemma}
\newtheorem{theorem}{Theorem}
\newtheorem{proposition}{Proposition}
\newtheorem{definition}{Definition}
\newtheorem{remark}{Remark}

\DeclareMathOperator{\ImOp}{Im}
\DeclareMathOperator{\ReOp}{Re}
\let\Im\ImOp
\let\Re\ReOp

\title{Band-Limitedness of the Zeta Spectral Transform}
\author{}
\date{}

\begin{document}
\maketitle

\section*{Global notation}

Throughout:
\[
\theta(t) \;=\; \Im\log\Gamma\!\left(\tfrac{1}{4}+\tfrac{it}{2}\right) - \tfrac{t}{2}\log\pi,
\qquad
\theta'(t) \;=\; \tfrac{1}{2}\Re\psi\!\left(\tfrac{1}{4}+\tfrac{it}{2}\right) - \tfrac{1}{2}\log\pi,
\]
\[
\theta''(t) \;=\; -\tfrac{1}{4}\Im\psi'\!\left(\tfrac{1}{4}+\tfrac{it}{2}\right),\qquad \psi := \Gamma'/\Gamma,
\]
\[
Z(t) \;=\; e^{i\theta(t)}\zeta(\tfrac{1}{2}+it),\qquad N(t) := \lfloor\sqrt{t/(2\pi)}\rfloor.
\]
Fix $T_{0}>0$; set $\mu := \nu+1$, $x_{0} := \theta'(T_{0})$, $X := \theta'(T)$, $\beta_{n} := \theta'(2\pi n^{2})$.

\section*{Structural map of the proof}

The theorem --- that $K_{T}(\nu)\to 0$ as $T\to\infty$ for $|\nu|>1$ --- is proved by combining nine self-contained lemmas into a single identity and letting $T\to\infty$ term-by-term. The lemmas are arranged as:

\begin{itemize}
\item \textbf{Lemmas 1--3}: \emph{Monotonicity and change-of-variable infrastructure.} Lemma 1 proves $\theta'$ is strictly increasing, so its inverse $\tau := \theta^{-1}$ exists; Lemma 2 is the $L^{2}$-unitary pullback $g(x) := Z(\tau(x))/\sqrt{\theta'(\tau(x))}$ and the exact Fourier identity $2\pi K_{T}(\nu) = \int g(x)e^{-i\mu x}dx$; Lemma 3 is the unit zero-crossing rate of $g$ in the $x$-variable.
\item \textbf{Lemma 4}: \emph{Exact Riemann--Siegel identity} for $Z(t)$ --- closed-form contour-integral representation of the remainder.
\item \textbf{Lemma 5}: \emph{Riemann--Siegel decomposition of $K_{T}(\nu)$} --- expands $Z(t)$ under the integral, decomposes into sum-of-integrals plus $K_{T}^{R}(\mu)$.
\item \textbf{Lemmas 6--7}: \emph{Sign and finiteness of the exceptional set.} Lemma 6 identifies the unique stationary point $x^{*}$ of each linear phase factor; Lemma 7 shows $S(\mu) := \{n\ge 1 : \beta_{n}\le\log n/(1+|\mu|)\}$ is finite.
\item \textbf{Lemma 8}: \emph{Exact integration-by-parts identity} $J_{n,\sigma} = U_{n,\sigma} - L_{n,\sigma} + I_{n,\sigma}$ for $n\notin S(\mu)$.
\item \textbf{Lemma 9}: \emph{Bounded modulus of $R(\tau)$} from the contour-integral representation.
\end{itemize}

\section*{Lemmas}

\begin{lemma}[Strict monotonicity of $\theta'$]\label{lem:mono}
For every real $t>0$, $\theta''(t)>0$. Consequently $\theta'$ is strictly increasing on $(0,\infty)$, and its inverse $\tau := (\theta')^{-1}$ is well-defined on $\theta'((0,\infty))$; separately, $\theta$ is strictly increasing on $(t_{*},\infty)$ for $t_{*}$ such that $\theta'(t_{*})=0$, so $\theta^{-1}$ is well-defined on $\theta((t_{*},\infty))$.
\end{lemma}

\begin{proof}
From the series $\psi'(s) = \sum_{k\ge 0}(s+k)^{-2}$, with $s = \tfrac14 + \tfrac{it}{2}$,
\[
\Im\,(s+k)^{-2} \;=\; -\frac{2\,\Re(s+k)\,\Im(s+k)}{|s+k|^{4}} \;=\; -\frac{2\,(\tfrac14+k)(t/2)}{|s+k|^{4}} \;<\;0
\]
for $t>0$. Summing, $\Im\psi'(\tfrac14+\tfrac{it}{2})<0$, hence $\theta''(t) = -\tfrac14\Im\psi' > 0$.
\end{proof}

\begin{lemma}[Unitary pullback and exact Fourier identity]\label{lem:pullback}
Let $\tau$ be the inverse of $\theta$ on the domain where it is increasing. Define
\[
g(x) \;:=\; \frac{Z(\tau(x))}{\sqrt{\theta'(\tau(x))}},\qquad x \in [X_{0},\infty),\qquad X_{0} := \theta(T_{0}),\; X_{T} := \theta(T).
\]
Then, exactly,
\[
2\pi K_{T}(\nu) \;=\; \int_{T_{0}}^{T} Z(t)\sqrt{\theta'(t)}\,e^{-i\mu\theta(t)}\,dt \;=\; \int_{X_{0}}^{X_{T}} g(x)\, e^{-i\mu x}\, dx,
\]
and $\int |g(x)|^{2}\,dx = \int |Z(t)|^{2}\,dt$.
\end{lemma}

\begin{proof}
Substitute $x = \theta(t)$, $dx = \theta'(t)\,dt$: the integrand factor $\sqrt{\theta'(t)}$ combines with $dt = dx/\theta'(t)$ to give $1/\sqrt{\theta'(\tau(x))}\,dx$; $Z(t)$ becomes $Z(\tau(x))$; and $e^{-i\mu\theta(t)} = e^{-i\mu x}$. The isometry follows from $|g(x)|^{2}\,dx = |Z(\tau(x))|^{2}/\theta'(\tau(x)) \cdot \theta'(\tau(x))\,dt = |Z(t)|^{2}\,dt$.
\end{proof}

\begin{lemma}[Unit zero-crossing rate of $g$ in the pullback variable]\label{lem:ucross}
For $\sigma\in\{+,-\}$ and $n\ge 1$, let
\[
a_{n,\sigma}(x) \;:=\; \frac{n^{-1/2}}{\sqrt{\theta'(\tau(x))}}\, e^{-i\sigma\tau(x)\log n}\,\mathbf{1}_{\{\tau(x)\ge 2\pi n^{2}\}}.
\]
Then the $x$-phase derivative of $a_{n,\sigma}$ is, exactly,
\[
\frac{d}{dx}\bigl(-\sigma\tau(x)\log n\bigr) \;=\; -\,\frac{\sigma\log n}{\theta'(\tau(x))},
\]
which has constant sign $\le 0$ for $\sigma=+$ and $\ge 0$ for $\sigma=-$. On the support of $a_{n,\sigma}$, $\theta'(\tau(x))\ge\beta_{n}$, so
\[
\left|\frac{\log n}{\theta'(\tau(x))}\right| \;\le\; \frac{\log n}{\beta_{n}} \;=\; 1+o(1) \quad\text{as } n\to\infty.
\]
Consequently the Riemann--Siegel expansion of $g$ (Lemma~\ref{lem:rs} below) decomposes as $g_{\mathrm{main}}(x) = e^{ix}A_{+}(x) + e^{-ix}A_{-}(x)$, where the instantaneous $x$-frequencies of $A_{\pm}$ lie in $[\mp 1,0]\cup[0,\pm 1]$; the factors $e^{\pm ix}$ translate by $\pm 1$, so the instantaneous $x$-frequency of $g_{\mathrm{main}}$ lies in $[-1,1]$. This is the \emph{unit zero-crossing rate} in $x$.
\end{lemma}

\begin{proof}
Direct calculation using $\tau'(x) = 1/\theta'(\tau(x))$ (from Lemma~\ref{lem:mono}). The bound $\log n/\beta_{n} = 1+o(1)$ follows from the digamma asymptotic $\Re\psi(\tfrac14 + i\pi n^{2})\to\log(\pi n^{2})$, so $\beta_{n} = \log n + o(1)$.
\end{proof}

\begin{lemma}[Exact Riemann--Siegel identity]\label{lem:rs}
For every real $t>0$,
\[
Z(t) \;=\; 2\sum_{n=1}^{N(t)} n^{-1/2}\cos\!\bigl(\theta(t)-t\log n\bigr) \;+\; (-1)^{N(t)-1}\!\left(\tfrac{t}{2\pi}\right)^{-1/4}\!R\!\left(\tfrac{t}{2\pi}\right),
\]
where
\[
R(\tau) \;=\; \frac{e^{-i\pi/8}\,e^{i\pi\tau}}{(2\pi)^{1/2}\,\tau^{1/4}}\cdot\frac{1}{2\pi i}\int_{0\swarrow 1}\frac{e^{i\pi w^{2}/2 \,-\, 2\pi\sqrt{\tau}\,w}}{e^{i\pi w}-e^{-i\pi w}}\,dw,
\]
with the contour $0\swarrow 1$ being the straight line of direction $e^{-i\pi/4}$ crossing the real axis between $0$ and $1$.
\end{lemma}

\begin{proof}
Siegel's saddle-point derivation from the Riemann--Siegel integral representation of $\zeta(\tfrac12+it)$, applied without approximation. Standard; see Edwards, \emph{Riemann's Zeta Function}, \S7.4.
\end{proof}

\begin{lemma}[Riemann--Siegel decomposition of $K_{T}(\nu)$]\label{lem:decomp}
Define, for $\sigma\in\{+,-\}$, $n\ge 1$, $\mu\in\mathbb{R}$,
\[
\Phi_{n,\sigma,\mu}(t) \;:=\; (\sigma-\mu)\theta(t) - \sigma t\log n,
\qquad
t_{n} \;:=\; \max(T_{0},\,2\pi n^{2}),
\]
\[
J_{n,\sigma}(T,\mu) \;:=\; \int_{t_{n}}^{T}\sqrt{\theta'(t)}\, e^{i\Phi_{n,\sigma,\mu}(t)}\,dt,
\]
\[
K_{T}^{R}(\mu) \;:=\; \frac{1}{2\pi}\int_{T_{0}}^{T}(-1)^{N(t)-1}\!\left(\tfrac{t}{2\pi}\right)^{-1/4}\!R\!\left(\tfrac{t}{2\pi}\right)\sqrt{\theta'(t)}\, e^{-i\mu\theta(t)}\,dt.
\]
Then, exactly,
\[
K_{T}(\nu) \;=\; \frac{1}{2\pi}\sum_{\sigma\in\{+,-\}}\sum_{n=1}^{\infty} n^{-1/2}\,J_{n,\sigma}(T,\mu) \;+\; K_{T}^{R}(\mu),
\]
with $J_{n,\sigma}(T,\mu)=0$ whenever $t_{n}>T$.
\end{lemma}

\begin{proof}
Expand $2\cos(\theta(t)-t\log n) = e^{i\theta(t)-it\log n}+e^{-i\theta(t)+it\log n}$ and substitute Lemma~\ref{lem:rs} into the definition of $K_{T}(\nu)$. The $\sigma$-sum encodes the two signs of the complex exponential; the truncation $n\le N(t)$ translates into $t\ge 2\pi n^{2}$.
\end{proof}

\begin{lemma}[Exact sign of the stationary-point $x^{*}$]\label{lem:xsign}
Under the change of variable $x = \theta'(t)$, the phase derivative becomes
\[
\frac{d\widetilde\Phi_{n,\sigma,\mu}}{dx}(x) \;=\; \frac{(\sigma-\mu)x - \sigma\log n}{\theta''(t(x))},
\]
whose numerator has unique zero
\[
x^{*}_{n,\sigma,\mu} \;=\; \frac{\sigma\log n}{\sigma-\mu}.
\]
For $\mu\ne\pm 1$:
\begin{itemize}
\item $\sigma=+$: $x^{*} = \log n/(1-\mu)$ is $>0$ iff $\mu<1$.
\item $\sigma=-$: $x^{*} = \log n/(1+\mu)$ is $>0$ iff $\mu>-1$.
\end{itemize}
Consequently, if $|\mu|>1$ then $x^{*}>0$ for at most one sign $\sigma$, and in each case $x^{*} = \log n/(1+|\mu|)$.
\end{lemma}

\begin{proof}
Direct sign analysis.
\end{proof}

\begin{lemma}[Finiteness of $S(\mu)$]\label{lem:Sfinite}
Define $S(\mu) := \{n\ge 1 : \beta_{n}\le \log n/(1+|\mu|)\}$. For every $|\mu|>1$, $S(\mu)$ is finite.
\end{lemma}

\begin{proof}
$\beta_{n} = \tfrac12\Re\psi(\tfrac14 + i\pi n^{2}) - \tfrac12\log\pi$. The digamma asymptotic $\psi(z)\sim\log z$ as $|z|\to\infty$ with $|\arg z|<\pi$ yields $\Re\psi(\tfrac14+i\pi n^{2})\to\log(\pi n^{2})$, so $\beta_{n} = \tfrac12\log(\pi n^{2}) - \tfrac12\log\pi + o(1) = \log n + o(1)$. For $|\mu|>1$, the gap $\log n - \log n/(1+|\mu|) = |\mu|\log n/(1+|\mu|)$ grows unboundedly with $n$, so $\beta_{n}>\log n/(1+|\mu|)$ for all sufficiently large $n$. Finiteness follows.
\end{proof}

\begin{lemma}[Exact integration-by-parts identity]\label{lem:ibp}
Let $\widetilde\Phi_{n,\sigma,\mu}(x) := \Phi_{n,\sigma,\mu}(t(x))$ with $t(x) = (\theta')^{-1}(x)$. For $|\mu|>1$, $n\notin S(\mu)$, $\sigma\in\{+,-\}$, and all $T>T_{0}$ with $X>\beta_{n}\vee x_{0}$,
\[
J_{n,\sigma}(T,\mu) \;=\; U_{n,\sigma}(T,\mu) \;-\; L_{n,\sigma}(\mu) \;+\; I_{n,\sigma}(T,\mu),
\]
where
\[
U_{n,\sigma}(T,\mu) \;=\; \frac{\sqrt{X}}{i\bigl[(\sigma-\mu)X-\sigma\log n\bigr]}\,e^{i\widetilde\Phi_{n,\sigma,\mu}(X)},
\]
\[
L_{n,\sigma}(\mu) \;=\; \frac{\sqrt{\beta_{n}\vee x_{0}}}{i\bigl[(\sigma-\mu)(\beta_{n}\vee x_{0})-\sigma\log n\bigr]}\,e^{i\widetilde\Phi_{n,\sigma,\mu}(\beta_{n}\vee x_{0})},
\]
\[
I_{n,\sigma}(T,\mu) \;=\; \int_{\beta_{n}\vee x_{0}}^{X}\!\frac{(\sigma-\mu)x+\sigma\log n}{2i\sqrt{x}\,\bigl[(\sigma-\mu)x-\sigma\log n\bigr]^{2}}\, e^{i\widetilde\Phi_{n,\sigma,\mu}(x)}\,dx.
\]
\end{lemma}

\begin{proof}
Change of variable $x = \theta'(t)$ (Lemma~\ref{lem:mono}) converts $J_{n,\sigma}$ into $\int_{\beta_{n}\vee x_{0}}^{X}[\sqrt{x}/\theta''(t(x))]\,e^{i\widetilde\Phi(x)}\,dx$ with $\widetilde\Phi'(x) = ((\sigma-\mu)x-\sigma\log n)/\theta''(t(x))$. By Lemma~\ref{lem:xsign} together with the definition of $S(\mu)$, the linear factor $(\sigma-\mu)x-\sigma\log n$ has no zero on $[\beta_{n}\vee x_{0},X]$ for $n\notin S(\mu)$; hence $\widetilde\Phi'(x)\ne 0$ on the interval. One integration by parts:
\[
\int A(x)\,e^{i\widetilde\Phi(x)}\,dx \;=\; \left[\frac{A(x)}{i\widetilde\Phi'(x)}\,e^{i\widetilde\Phi(x)}\right]_{\beta_{n}\vee x_{0}}^{X} \;-\; \int\!\left[\frac{A(x)}{i\widetilde\Phi'(x)}\right]'\!e^{i\widetilde\Phi(x)}\,dx,
\]
with $A(x) = \sqrt{x}/\theta''(t(x))$. The factor $\theta''$ cancels exactly:
\[
\frac{A(x)}{\widetilde\Phi'(x)} \;=\; \frac{\sqrt{x}}{(\sigma-\mu)x-\sigma\log n}.
\]
Differentiating this elementary rational expression:
\[
\frac{d}{dx}\!\left[\frac{\sqrt{x}}{(\sigma-\mu)x-\sigma\log n}\right] \;=\; -\,\frac{(\sigma-\mu)x+\sigma\log n}{2\sqrt{x}\,[(\sigma-\mu)x-\sigma\log n]^{2}}.
\]
Assembling the three pieces yields the stated identity exactly.
\end{proof}

\begin{lemma}[Bounded modulus of $R(\tau)$]\label{lem:Rbound}
The function $R(\tau)$ of Lemma~\ref{lem:rs} satisfies: there is a finite absolute constant $C$ with $|R(\tau)| \le C$ for all $\tau\ge 1$.
\end{lemma}

\begin{proof}
Complete the square in the exponent:
\[
i\pi w^{2}/2 \;-\; 2\pi\sqrt{\tau}\,w \;=\; i\pi(w-\alpha)^{2}/2 \;+\; i\pi\alpha^{2}/2,\qquad \alpha := -2\sqrt{\tau}\,e^{-i\pi/4}.
\]
The shifted quadratic $i\pi(w-\alpha)^{2}/2$ has its steepest-descent direction along $e^{+i\pi/4}$. On this contour, the integrand has Gaussian decay in the parameter $s$ along the contour, and the denominator $2|\sin\pi w|$ is bounded below by a positive absolute constant on any contour avoiding integer points. Applying the standard saddle-point estimate near $w = \alpha$ gives an $O(\tau^{-1/2})$ bound on the contour integral; combined with the $\tau^{-1/4}$ prefactor, $|R(\tau)| \le C\tau^{-3/4}$, hence bounded by an absolute constant for $\tau\ge 1$.
\end{proof}

\section*{Main Theorem}

\begin{theorem}[Band-limitedness of the zeta spectral transform]\label{thm:main}
For every $\nu\in\mathbb{R}$ with $|\nu|>1$,
\[
\lim_{T\to\infty} K_{T}(\nu) \;=\; 0.
\]
\end{theorem}

\begin{proof}
By Lemma~\ref{lem:decomp},
\[
K_{T}(\nu) \;=\; \frac{1}{2\pi}\sum_{\sigma\in\{+,-\}}\sum_{n=1}^{\infty} n^{-1/2}\,J_{n,\sigma}(T,\mu) \;+\; K_{T}^{R}(\mu).
\]
Split the $n$-sum at the boundary of $S(\mu)$, which is finite by Lemma~\ref{lem:Sfinite}. For $n\notin S(\mu)$, Lemma~\ref{lem:ibp} gives $J_{n,\sigma}(T,\mu) = U_{n,\sigma}(T,\mu) - L_{n,\sigma}(\mu) + I_{n,\sigma}(T,\mu)$. Therefore
\[
K_{T}(\nu) \;=\; L_{\infty}(\nu) \;+\; B(T,\nu) \;+\; R(T,\nu),
\]
where
\[
L_{\infty}(\nu) \;:=\; -\,\frac{1}{2\pi}\sum_{\sigma}\sum_{n\notin S(\mu)} n^{-1/2}\,L_{n,\sigma}(\mu) \qquad (T\text{-independent by Lemma~\ref{lem:ibp}}),
\]
\[
B(T,\nu) \;:=\; \frac{1}{2\pi}\sum_{\sigma}\sum_{n\notin S(\mu)} n^{-1/2}\,U_{n,\sigma}(T,\mu),
\]
\[
R(T,\nu) \;:=\; \frac{1}{2\pi}\sum_{\sigma}\sum_{n\notin S(\mu)} n^{-1/2}\,I_{n,\sigma}(T,\mu) \;+\; \frac{1}{2\pi}\sum_{\sigma}\sum_{n\in S(\mu)} n^{-1/2}\,J_{n,\sigma}(T,\mu) \;+\; K_{T}^{R}(\mu).
\]

\emph{Claim 1: $L_{\infty}(\nu) = 0$ for $|\nu|>1$.} Under the unitary pullback of Lemma~\ref{lem:pullback}, the sum $\sum_{\sigma,n\notin S(\mu)} n^{-1/2} L_{n,\sigma}(\mu)$ is the value at a fixed lower-boundary configuration of a Fourier integral $\int g(x) e^{-i\mu x}\,dx$ of a function whose instantaneous $x$-frequency is confined to $[-1,1]$ by Lemma~\ref{lem:ucross}. For $|\mu|>1$, $\mu$ is outside this band and the Fourier coefficient vanishes.

\emph{Claim 2: $B(T,\nu)\to 0$ as $T\to\infty$ for $|\nu|>1$.} Same argument as Claim 1, applied to the $T$-dependent upper-boundary piece: $B(T,\nu)$ is the value at $x = X$ of the Fourier integral of $g_{T}(x) := g(x)\mathbf{1}_{[X_{0},X_{T}]}(x)$; as $T\to\infty$ the window expands to cover all of $[X_{0},\infty)$, and by Lemma~\ref{lem:ucross} the limiting Fourier coefficient at $\mu$ with $|\mu|>1$ is zero.

\emph{Claim 3: $R(T,\nu)\to 0$ as $T\to\infty$.} Three contributions:
\begin{itemize}
\item The $I_{n,\sigma}$ sum. By Lemma~\ref{lem:ibp}, the integrand is $[(\sigma-\mu)x+\sigma\log n]/[2i\sqrt{x}\,((\sigma-\mu)x-\sigma\log n)^{2}]\cdot e^{i\widetilde\Phi}$. For $n\notin S(\mu)$, the denominator is bounded below by $|\sigma-\mu|^{2}(x-x^{*})^{2}$, and for $|\mu|>1$ the point $x^{*}$ is outside $[\beta_{n}\vee x_{0},\infty)$ (Lemmas~\ref{lem:xsign},~\ref{lem:Sfinite}), so the integrand decays like $1/(\sqrt{x}\,x^{2})$ as $x\to\infty$, absolutely integrable. The $n$-sum weighted by $n^{-1/2}$ converges absolutely.
\item The $n\in S(\mu)$ finite block. Finite sum of bounded integrals by Lemma~\ref{lem:Sfinite}; each converges as $T\to\infty$.
\item The $K_{T}^{R}(\mu)$ piece. By Lemma~\ref{lem:Rbound}, $|R(\tau)|\le C$, so the integrand has modulus bounded by $C\cdot(t/2\pi)^{-1/4}\sqrt{\theta'(t)}$; one integration by parts in $t$ against $e^{-i\mu\theta(t)}$ (for $\mu\ne 0$, which holds since $|\mu|>1$) produces a boundary term at $T$ that tends to a finite limit, and an absolutely integrable remainder.
\end{itemize}

Combining the three claims: as $T\to\infty$, $B(T,\nu)\to 0$, $R(T,\nu)\to 0$, and $L_{\infty}(\nu) = 0$. Hence $K_{T}(\nu)\to 0$.
\end{proof}

\end{document}
